\cotefile{156-9}
\scanpage{119}
\mspage{113}
% archive sheet number: 118

\uline{Lemme 2} Avec les notations précédentes, les conditions suivantes sont équivalentes

(i) $A\cup B\subset\Sigma'$, où $\Sigma'$ est une partie de $\Sigma$ stable par $\operatorname{Sup}$, $\operatorname{Inf}$ et $\unsure{\ell|\complement}$, et qui est une alg. de Boole pour l'ordre induit par $\Sigma$ et pour l'op. induite par $\ell$.

(ii) On a

 a) $\forall\deleted{\unsure{b}}a\in\deleted{\unsure{B}}A$, on a $a=a'\vee a''$

 b) \deleted{La famille \unsure{formée} des $b\in\unsure{I|B}$} L'ensemble \deleted{\unsure{formé}}\deleted{\unsure{partie}}\deleted{de $B$}
\[
\unsure{\widehat B|\widetilde B}=B\cup\{a''\mid a\in A\text{ t.q. }a''\ne0\}
\]
est admissible.

Sous ces conditions, l'alg. de Boole $\Sigma_{\widehat B}$ engendrée par $\widehat B$ (cf déf. p.\unsure{110|140}) \deleted{est aussi l'} \unsure{iso.} à $\unsure{\mathcal P|\mathscr P}(\widehat B)$, est aussi le plus petit des sous-ens. $\Sigma'$ satisfaisant (i).

\margin{\uline{Déf}}
On dit alors que $\{A,B\}$ \deleted{\unsure{sont}}\unsure{est} admissible.

\uline{Axiome \unsure{Atm. g)|Atm. 9)} $-\unsure{\gamma\gamma^2|\gamma^2}$} Pour toute \deleted{$\unsure{F}$} $\unsure{\mathbb G|G}\ll\unsure{\mathbb F|F}$ en $F\unsure{,}G\in\unsure{\mathcal F|\mathscr F}$, posant
\[
B=\{\unsure{X^0|\mathbb X^0|X^\circ|\mathbb X^\circ}\mid\unsure{X|\mathbb X}\in G\},\quad A=\{X^0\mid X\in F\},
\]
